
\chapter{Exercise Specification}
\label{app:exercise-spec}

RekenRangers supports three arithmetic operations: addition, subtraction, and multiplication. For each operation, exercises are classified into three difficulty levels based on properties of the operands and the result. This appendix describes the classification rules for each operation.

\section{Addition}

The difficulty of an addition exercise $a + b = s$ is determined by the magnitude of the operands and whether the addition requires carrying (i.e., whether the sum of the unit digits exceeds or equals 10).

\begin{table}[ht]
\centering
\begin{tabular}{cp{10cm}}
\toprule
Level & Criteria \\
\midrule
1 (easy) & The sum is less than 10. Both operands are single-digit numbers and the result stays within the single digits. \\
2 (medium) & The sum is 10 or greater, but no carrying is required, except for two single digits. This includes cases where both operands are single-digit and the sum just crosses 10, as well as cases involving two-digit operands where the unit digits sum to less than 10. \\
3 (difficult) & Both operands are two-digit numbers and carrying is required: the sum of the unit digits is 10 or greater. \\
\bottomrule
\end{tabular}
\caption{Difficulty classification for addition exercises.}
\label{tab:diff-addition}
\end{table}

\section{Subtraction}

The difficulty of a subtraction exercise $a - b$ (where $a \geq b$) depends on the magnitude of the operands and whether borrowing is required (i.e., whether the unit digit of $a$ is smaller than the unit digit of $b$).

\begin{table}[ht]
\centering
\begin{tabular}{cp{10cm}}
\toprule
Level & Criteria \\
\midrule
1 (easy) & Both operands are single-digit numbers and the result is less than 10. \\
2 (medium) & At least one operand is a two-digit number, but no borrowing is required: the unit digit of $a$ is greater than or equal to the unit digit of $b$. \\
3 (difficult) & Both operands are two-digit numbers and borrowing is required: the unit digit of $a$ is smaller than the unit digit of $b$. \\
\bottomrule
\end{tabular}
\caption{Difficulty classification for subtraction exercises.}
\label{tab:diff-subtraction}
\end{table}

\section{Multiplication}

The difficulty of a multiplication exercise $a \times b = p$ is based on the magnitude of the operands and the product.

\begin{table}[ht]
\centering
\begin{tabular}{cp{10cm}}
\toprule
Level & Criteria \\
\midrule
1 (easy) & The product is at most 25. These are small multiplication facts that children typically memorise early. \\
2 (medium) & Both operands are at most 10 (single-digit multiplication tables), but the product exceeds 25. \\
3 (difficult) & Both operands are at most 20, with exactly one operand exceeding 10. These exercises require multiplying a single-digit number by a two-digit number. \\
\bottomrule
\end{tabular}
\caption{Difficulty classification for multiplication exercises.}
\label{tab:diff-multiplication}
\end{table}

\noindent These classification rules are implemented as deterministic functions in the back-end. When the adaptive system selects an exercise, it first determines which difficulty level to draw from (based on the child's current rating and the distribution in Table~\ref{tab:difficulty-zones}), then generates a random exercise that satisfies the corresponding criteria.\\ \\
The difficulty classification described above is based on the categorisation of arithmetic exercises used in the Flemish primary school mathematics curriculum \cite{REF_CURRICULUM}, which distinguishes exercises by operand magnitude and the presence or absence of carrying and borrowing as key indicators of cognitive complexity.
